\chapter{Introduction}\label{ch:introduction}
Some introducing words \dots


\section{Motivation}\label{sec:motivation}

\subsection{The Agentic Shift}\label{subsec:the-agentic-shift}
- **From Chat to Action:** The evolution of LLMs from static information retrieval to goal-oriented autonomous agents executing multi-step workflows.
- **The Trajectory:** Defining the unit of analysis not as an output string, but as a temporal sequence: $\tau = (Observation, Thought, Action, Outcome)$.

\subsection{The Evaluation Gap}\label{subsec:the-evaluation-gap}
- **Static vs. Dynamic:** Why fixed benchmarks (MMLU, GSM8K) fail to evaluate agentic properties like error recovery and tool usage.
- **The "Scalar Fallacy":** The inadequacy of single-number ELO ratings to capture orthogonal capabilities (e.g., a "Creative" agent vs. a "Compliant" agent) or intransitive "rock-paper-scissors" dynamics.

\section{Problem Statement}\label{sec:problem-statement}


\section{Research Goal}\label{sec:research-goal}
1. **System Design:** Develop a **Hybrid Semantic Router** to dispatch agents to domain-specific judges.
2. **Ranking Methodology:** Implement a **Vectorized ELO (OpenSkill)** engine to handle intransitive skill sets.
3. **Validation:** Demonstrate the framework's convergence efficiency using a **Procedurally Generated Agent Population** evaluated against a diverse stream of **User Queries**.

\section{Outline of Work}\label{sec:outline-of-work}
What is the structure of this thesis?